\chapter{Results}

\section{Overview}

Results were obtained from the derivative output directory
\texttt{/Volumes/OHDD/DATA/epilepsy/derivatives/seegring/sub-EP01AN96M1047}.
Two complementary result summaries were available. First, the batch
condition-discrimination report generated on 2026-03-24 summarized
population-level analyses across 9 cognitive tasks and 11 runs.
Second, a separate channel-significance summary quantified
channel-by-time cluster significance and decoding-based channel
importance for a subset of 7 analyzable task runs.

These two summaries support a consistent interpretation. At the
population latent-space level, no task showed statistically significant
condition separation after cluster-corrected permutation testing. In
contrast, at the channel level, several tasks exhibited substantial
numbers of channels with significant time clusters and robust decoding
performance, indicating that condition information was present in the
recordings but was not recovered as a statistically significant
population-level latent separation under the current analysis settings.

\section{Population-Level Condition Separation}

The cross-task report in
\texttt{reports/20260324\_171402/report.md} analyzed 9 cognitive tasks
spanning 11 runs, including lexical decision, shape control, sentence
noun, sentence grammar, salience pain, balloon watching, viseme
perception, and viseme generation. Statistical significance was assessed
with cluster-based permutation tests using 1000 permutations and
$\alpha=0.05$.

The main result of that report was negative: no task showed
statistically significant condition-dependent separation at the
population level. The corresponding summary table
(\texttt{stats\_summary.csv}) reported \texttt{Sig = NO} for every
processed task, with peak separation, minimum p-value, peak decoding,
decoding onset, effective dimensionality, and jPCA summary fields left
as undefined values. The report therefore does not support any claim of
significant condition separation in latent space, significant
condition-dependent temporal windows, or significant cross-condition
population decoding at the whole-task summary level.

\section{Channel-Level Significant Results}

A separate analysis stored in
\texttt{channel\_significance\_summary.mat} yielded a different picture
at the channel level. This analysis re-ran permutation statistics and
decoding for 7 task runs and computed the number of channels with
significant channel-by-time clusters. All 7 analyzed tasks showed at
least some cluster-corrected channel-level significance, although none
showed region-of-interest significance after false-discovery-rate
correction.

The strongest channel-level effects were observed for
\texttt{task-lexicaldecision\_run-01}, which showed 23 channels with
significant time clusters and a peak decoding accuracy of 92.0\%.
\texttt{task-shapecontrol\_run-01} showed a similarly strong effect,
with 22 significant channels and a peak decoding accuracy of 99.0\%.
Intermediate effects were observed for
\texttt{task-sentencegrammar\_run-01} and
\texttt{task-viseme\_run-01}, each with 9 significant channels, and for
\texttt{task-saliencepain\_run-01} and
\texttt{task-sentencenoun\_run-01}, each with 7 significant channels.
\texttt{task-balloonwatching\_run-01} showed 5 significant channels and
a nominal peak decoding accuracy of 100.0\%, although this result should
be interpreted cautiously given the very small sample size of 10 trials.

Across these 7 channel-level analyses, the count of ROI-level FDR
significant channels was zero in every task. Thus, the significant
effects reported here should be interpreted specifically as
channel-by-time cluster-corrected findings, not as ROI-level
FDR-significant contacts. The effects were spatially sparse and
time-localized rather than strong enough to survive ROI-level
multiple-comparison correction.

\section{Decoding Results}

The channel-level decoding analysis indicated that binary tasks carried
considerable condition information despite the absence of significant
population-level latent separation in the batch report. Peak decoding
accuracy reached 99.0\% in shape control, 92.0\% in lexical decision,
75.0\% in salience pain, 68.3\% in sentence grammar, and 65.0\% in
sentence noun. Balloon watching reached 100.0\%, but this estimate was
derived from only 10 trials and is therefore likely to be unstable.

By contrast, the multi-class viseme task showed much weaker decoding,
with a peak accuracy of 16.4\% across 14 classes. Although this is above
the nominal chance level of approximately 7.1\%, it was accompanied by
only modest channel-level clustering and was not reflected in a
significant cross-task population-level result.

\section{Spatial Distribution of Informative Channels}

The most repeatedly prominent channels across tasks were concentrated in
a relatively stable subset of contacts. Among the top-20 channels ranked
by decoding contribution, \texttt{PI6}, \texttt{PI7}, \texttt{PI5},
\texttt{PI4}, \texttt{Ia15}, and \texttt{PI3} each appeared in 6 of the
7 analyzed tasks. Additional frequently recurring channels included
\texttt{Ia14} in 5 tasks and \texttt{PI8}, \texttt{PI9}, \texttt{C14},
\texttt{C13}, \texttt{C12}, \texttt{Ia16}, and \texttt{PI2} in 4 tasks.

This recurrence suggests that a stable subset of electrodes carried much
of the discriminative signal across paradigms. The recurrent involvement
of neighboring \texttt{PI}, \texttt{Ia}, and \texttt{C} contacts is more
consistent with localized informative subspaces than with a broadly
distributed whole-population separation effect.

\section{Interpretation}

Taken together, the current outputs support a split conclusion. On the
one hand, the official batch condition-discrimination report provides no
evidence for statistically significant condition separation at the
population latent-space level after cluster-corrected testing. On the
other hand, the dedicated channel-significance analysis shows that
multiple tasks do contain statistically significant channel-by-time
effects and practically strong decoding, especially in lexical decision
and shape control.

The most conservative interpretation is therefore that condition-related
information is detectable in localized channels and in decoding space,
but that this information was not sufficiently strong, stable, or
distributed to yield significant whole-population latent separation in
the current batch summary. This distinction should be preserved in the
manuscript: significant channel-level effects are supported by the saved
outputs, but those channel-level effects do not survive the stricter
ROI-level FDR criterion, whereas significant population-level condition
separation is also not supported by the saved batch report.
